\paragraph{Theorem E.2.9 (Proved: framed null-bordism over the fundamental group).}
Let $G$ be the group in (E.16). There is a closed smooth
aspherical four-manifold $N$ with
\[
 \pi_1(N)=G,\qquad \widetilde N\simeq\mathbb H^2\times\mathbb R^2,
 \qquad H_*(N;\mathbb Z)\cong H_*(S^1\times S^3;\mathbb Z),
 \tag{E.B3}
\]
and a smooth classifying map $c:P\to N$ inducing the displayed
identification of fundamental groups. With the natural stable
framing of E.2.8, one has
\[
                 [P,c]=0\quad\hbox{in }\Omega^{\mathrm{fr}}_{10}(N).
 \tag{E.B4}
\]
Thus $P$ has a compact stably framed eleven-dimensional filling
over $BG$. This is stronger than the filling of E.2.8; it is
still a bordism statement, not a diffeomorphism classification.
\begin{proof}
\emph{The aspherical target.}
The repaired quotient maps of A.2.1 are smooth maps
$e_i:S^6\to S^3$ over the base, and B.2.1 gives the exact
factorizations $f_i=h_i e_i$. Here $h_i:S^3\to\mathbb P^1$
is the weighted Hopf map of orders $3,4$, with the second
map shifted by the chosen M\"obius transformation. Define
\[
 N=\{(u,v)\in S^3\times S^3:h_1(u)=h_2(v)\},\qquad
 c(x,y)=(e_1(x),e_2(y)).
 \tag{E.B5}
\]
The critical values of each $h_i$ are only its two exceptional
circle values. They are disjoint, so the same transversality
argument as in E.2.2 makes $N$ a compact smooth four-manifold.
The two circle actions on the weighted Hopf spheres give a
principal $T^2$ bundle in the orbifold sense over
\[
                  \mathcal B=S^2(3,4,3,4).
\]
At a cone of order $m$, its uniformizing chart is
$(D^2\times T^2)/C_m$, with the generator rotating the disc
and translating one torus coordinate by a primitive fraction
of denominator $m$. This local description is smooth and is
the same finite-filling description as A.1.3, with the unused
circle retained.

The orbifold has Euler characteristic
$2-2(1-1/3)-2(1-1/4)=-5/6$. It is good and has hyperbolic
orbifold universal cover, by
\cite[Theorems 2.3--2.4 and the Euler-number discussion in
Section 2]{rootScott83}. Pulling the principal torus bundle to
$\mathbb H^2$ gives an ordinary principal $T^2$ bundle over a
contractible base, hence the trivial bundle. Its universal
cover is $\mathbb H^2\times\mathbb R^2$. The local chart above
shows that these orbifold pullbacks are ordinary covering maps
of total spaces. This proves asphericity and injectivity of
the fibre subgroup $\mathbb Z^2$.

Using the marked fibre generators $h_1,h_2$ and the four
positive peripheral generators gives exactly
\[
 \begin{split}
 \pi_1(N)=\langle a_1,a_2,b_1,b_2,h_1,h_2\mid {}&h_1,h_2
       \text{ central},\quad a_1^3=h_1,\quad a_2^4=h_1^{-1},\\
 &b_1^3=h_2,\quad b_2^4=h_2^{-1},\quad
 a_1a_2b_1b_2=1\rangle.
 \end{split}
\]
The map $c$ sends the surviving regular-fibre generators
$\widehat\gamma,\widehat\gamma'$ to $h_1,h_2$, respectively,
and each peripheral generator to the one with the same name.
This is because $e_i$ is the literal circle quotient on these
fibres and finite collars. Its cusp replacement extends over
a solid torus and kills precisely the cusp meridian in this
presentation. Comparing with the full van Kampen calculation
in E.2.3 therefore shows that $c_*$ is the identity on the
displayed presentation, and in particular is an isomorphism.
As $N$ is aspherical, $c$ is a classifying map.

For completeness the abelian relation matrix has Smith
diagonal $(1,1,1,1,1)$; a primitive integer kernel vector in
the order $(a_1,a_2,b_1,b_2,h_1,h_2)$ is
$(4,-3,-4,3,12,-12)$. Thus $H_1(N)=\mathbb Z$.
Every stratum of its base has a torus fibre, so Euler
additivity gives $\chi(N)=0$. Integral Poincare duality and
the universal coefficient theorem give $H_3(N)=\mathbb Z$
and no torsion in $H_2(N)$. Euler characteristic then forces
$H_2(N)=0$, proving the integral homology assertion in (E.B3).

\emph{The bordism computation before taking the inverse image.}
For a based sphere, the stable splitting
$(S^3)_+\simeq S^0\vee S^3$ identifies
\[
 \Omega^{\mathrm{fr}}_6(S^3)
       =\pi^{\mathrm S}_6\oplus\pi^{\mathrm S}_3.
 \tag{E.B6}
\]
We use Pontryagin--Thom in its form for maps to a space,
\cite[Definition 8.16 and Theorem 8.17]{rootDavisKirk01}.
The first component of $[S^6,e_i]$ is obtained by forgetting
the map, so it is zero for the bounding sphere framing.
The second component is the stable class of $e_i$, equivalently
the framed inverse image of a regular point. It is $\eta^3$
by B.1.2. One can also check the splitting directly: the
Pontryagin--Thom collapse of the framed sphere into its
disjoint-basepoint suspension has degree one on the sphere
summand; composing with $e_i$ gives its stable homotopy class.
Thus
\[
                         [S^6,e_i]=(0,\eta^3).
\]
Smashing the two sphere splittings gives
\[
 (S^3\times S^3)_+\simeq
       S^0\vee S^3\vee S^3\vee S^6.
\]
Under this splitting the external product
$[S^6\times S^6,e_1\times e_2]$ has four components.
The first three each contain a factor $[S^6]=0$; the last
is $\eta^3\eta^3=\eta^6=0$. Indeed $\eta^4$ belongs to
$\pi^{\mathrm S}_4=0$, as recorded in
\cite[Theorem 2.1 and Table 1]{rootIWX20}. Therefore
\[
 [S^6\times S^6,e_1\times e_2]=0
      \quad\hbox{in }\Omega^{\mathrm{fr}}_{12}(S^3\times S^3).
 \tag{E.B7}
\]
This calculation uses the product of the specific bounding
sphere framings, not just vanishing of ordinary characteristic
classes.

\emph{The comparison with the requested classifying map.}
By (E.B7) there is a compact stably framed thirteen-manifold
$Z$ with boundary $S^6\times S^6$ and a map
$E:Z\to S^3\times S^3$ extending $e_1\times e_2$.
Smooth it relative to the boundary and make it transverse
to $N$ there. Transversality on the boundary already holds:
the derivative of its defining equation is $df_1-df_2$,
which is onto. Relative transversality therefore supplies
the interior perturbation without changing that boundary.

The normal bundle of $N\subset S^3\times S^3$ is the
restriction of $h_1^*T\mathbb P^1$ pulled back from the first
sphere. This oriented two-plane bundle is trivial because
$H^2(S^3;\mathbb Z)=0$. Choose a framing on it before
restriction to $N$. Its pullback under $e_1\times e_2$
is a framing of $f_1^*T\mathbb P^1$, agreeing with that
in E.2.8 up to homotopy since $H^1(S^6;\mathbb Z)=0$.
It follows that
\[
                  V=E^{-1}(N),\qquad E|_V:V\longrightarrow N
\]
is a compact stably framed eleven-manifold with boundary
$(P,c)$ and the required boundary framing. This proves (E.B4).

The exact group-matrix and Euler calculation, together with
the two equivalent stable-product reductions, is recorded in
\texttt{code/shifted\_group\_bordism.py}.
The earlier disc-extension filling need not extend the
classifying map; (E.B7) supplies a different filling that does.
Neither construction computes the higher equivariant surgery
data needed to distinguish smooth ten-manifolds over $N$.
\end{proof}
